\documentclass[11pt]{article}

\usepackage[T1]{fontenc}
\usepackage[utf8]{inputenc}
\usepackage{amsmath,amssymb,amsthm}
\usepackage{array,booktabs}
\usepackage[margin=1in]{geometry}
\usepackage[hidelinks]{hyperref}

\newtheorem{theorem}{Theorem}
\newtheorem{lemma}[theorem]{Lemma}
\newtheorem{proposition}[theorem]{Proposition}
\newtheorem{corollary}[theorem]{Corollary}
\theoremstyle{definition}
\newtheorem{definition}[theorem]{Definition}
\newtheorem{remark}[theorem]{Remark}

\newcommand{\DTS}{\mathsf{DTS}}
\newcommand{\QBF}{\mathsf{QBF}}
\newcommand{\SAT}{\mathsf{SAT}}
\newcommand{\leanus}{\textunderscore\allowbreak}
\newcommand{\lean}[1]{\ifmmode\text{\ttfamily #1}\else{\ttfamily\let\_\leanus #1}\fi}

\title{The alternation-trading calculus for $\QBF_k$ versus $\DTS$\\ stops at Williams's exponent: a kernel-checked limit}
\author{PastaPastaPasta}
\date{November 2026}

\emergencystretch=2em

\begin{document}
\maketitle

\begin{abstract}
Williams \cite[Section~5]{W13} applies alternation-trading proofs to show that for every integer $k\ge1$, $\QBF_k\notin\DTS[n^c]$ whenever $c>k$ satisfies $c^3/k-c^2-2c+k<0$ (Theorem~5.1). For $k=1$ ($\SAT$ versus $\DTS[n^c]$), the unique root in $(1,2)$ is $r_1=2\cos(\pi/7)\approx1.801938$, which Buss and Williams \cite{BW15} proved is the exact limit of the alternation-trading calculus. For general $k\ge2$, Williams's exponents are $c_2\approx2.903212$, $c_3\approx3.942820$, $c_4\approx4.962389$, and so on, and no upper limit for the $k$-block calculus was established.

We prove in Lean~4 that for every integer $k\ge1$, a formalisation \lean{QStep} of the $\QBF_k$ alternation-trading calculus has no refutation for any $c>k$ with $c^3/k-c^2-2c+k\ge0$ (that is, $c\ge c_k$), and strictly above $c_k$ excludes even non-strict refutations ($a'\le a_0$ from $a_0>0$ in at least one step). The proof constructs a $k$-parametric grouping potential over partitions of the quantifier stack into consecutive pop groups of size $1,\dots,k$; at $k=1$ it recovers the Buss--Williams bound $2\cos(\pi/7)$ directly. On the lower side, we provide a rational transcription \lean{checkL} of Williams's literal rules, prove its soundness into \lean{QStep}, and kernel-check explicit refutation certificates at $c=7/4$ ($k=1$), $c=289/100$ ($k=2$), $c=393/100$ ($k=3$), and $c=989/200$ ($k=4$), bracketing the formal thresholds in $[1.75, c_1]$, $[2.89, c_2]$, $[3.93, c_3]$, and $[4.945, c_4]$.

\emph{What is kernel-checked:} the upper limit $c\ge c_k$ for all $k\ge1$ (and non-strict exclusion for $c>c_k$), the equality $c_1=2\cos(\pi/7)$, the soundness of the literal checker \lean{checkL} into \lean{QStep}, and the four literal refutation certificates.
\emph{What is not machine-checked:} that the formal rules coincide with Williams's brief prose description in \cite[Section~5]{W13} (a modelling claim set out in Section~\ref{sec:calculus}), Williams's own parametric family achieving every $c<c_k$, and the complexity-theoretic simulation lemmas underlying \cite{W13}. The limit applies only to this formal calculus.
\end{abstract}

\section{Introduction}

\paragraph{Background.}
In \cite{W13}, Williams formalised alternation-trading proofs for deterministic time--space lower bounds of the form $\SAT\notin\DTS[n^c]$ (time $n^{c+o(1)}$, space $n^{o(1)}$) and extended them in Section~5 to quantified Boolean formulas with $k$ alternating quantifier blocks ($\QBF_k$). Theorem~5.1 of \cite{W13} establishes $\QBF_k\notin\DTS[n^c]$ for every $c<c_k$, where $c_k$ is the unique root in $(k,k+1)$ of
\[
p_k(c) := \frac{c^3}{k} - c^2 - 2c + k = 0.
\]
Table~\ref{tab:roots} lists the roots $c_k$ and our kernel-checked certificate values for $k=1,2,3,4$. In the CMU full version of \cite{W13} (end of Section~3), Williams remarked empirically that automated search did not find a better lower bound than $c_k$, though Section~5 of \cite{W13} states no formal optimality theorem or conjecture. For $k=1$, Buss and Williams \cite{BW15} proved that $c_1=2\cos(\pi/7)$ is indeed the exact limit of the calculus.

\begin{table}[h]
\centering
\begin{tabular}{@{}ccccccc@{}}
\toprule
$k$ & Williams root $c_k$ & $q_k = k/c_k$ & Certified $c$ & Lines & Final $a'/a_0$ & Lean certificate \\
\midrule
1 & $1.801938$ & $0.554958$ & $7/4 = 1.75$ & 140 & $0.801670$ & \lean{qcert1\_ok} \\
2 & $2.903212$ & $0.688892$ & $289/100 = 2.89$ & 324 & $0.970106$ & \lean{qcert2\_ok} \\
3 & $3.942820$ & $0.760877$ & $393/100 = 3.93$ & 432 & $0.983568$ & \lean{qcert3\_ok} \\
4 & $4.962389$ & $0.806063$ & $989/200 = 4.945$ & 540 & $0.969004$ & \lean{qcert4\_ok} \\
\bottomrule
\end{tabular}
\caption{Williams's $\QBF_k$ exponents $c_k$ (roots of $c^3/k-c^2-2c+k=0$ in $(k,k+1)$) and our kernel-checked literal certificates.}\label{tab:roots}
\end{table}

\paragraph{Results.}
We determine the upper limit of the formal $\QBF_k$ alternation-trading calculus for every $k\ge1$ in Lean~4:
\begin{enumerate}
\item For every integer $k\ge1$ and every real $c>k$ with $c^3/k-c^2-2c+k\ge0$, no path in the relaxed calculus \lean{QStep} from $\DTS[n^{a_0}]$ ($a_0\ge0$) back to a quantifier-free line $\DTS[n^{a'}]$ achieves $a'<a_0$ (\lean{qbf\_no\_refutation}, \lean{qbf\_no\_refutation\_reach}, \lean{not\_refutes}).
\item When $c^3/k-c^2-2c+k>0$, even non-strict refutations ($a'\le a_0$ from $a_0>0$ in at least one step) are impossible (\lean{qbf\_no\_refutation\_strict}).
\item At $k=1$, the unique threshold parameter satisfies $c=2\cos(\pi/7)$ (\lean{par\_one\_eq\_r7}), recovering the Buss--Williams bound for $\SAT$ versus $\DTS$ (\lean{sat\_no\_refutation}).
\item For $k=1,2,3,4$, literal refutations checked by \lean{checkL} (tracking guesses, input floor $1$, and $x\le a-1$) verify the lower endpoints $c=7/4, 289/100, 393/100, 989/200$ via \lean{decide +kernel}, with soundness proved into \lean{QStep} (\lean{checkL\_sound}, \lean{interval\_k1}--\lean{interval\_k4}).
\end{enumerate}

\paragraph{AI assistance.}
The proofs and Lean formalisation were produced with AI assistance under the author's direction. The kernel-checked Lean proofs in \lean{PastaMath/QbfThreshold/} are the primary evidence for the theorems about the formal calculus.

\section{The calculus}\label{sec:calculus}

\subsection{Transcription from Williams (TOCT 2013, Section~5)}

Williams \cite[Section~5]{W13} describes the $\QBF_k$ calculus in one sentence: only the Slowdown Rule differs from the $\SAT$ calculus (Definition~4.1 of \cite{W13}), because its application removes up to $k$ quantifier blocks instead of just one.

A literal line is an alternating class
\[
(Q_1\,n^{a_1})^{b_2}(Q_2\,n^{a_2})^{b_3}\cdots(Q_m\,n^{a_m})^{b_{m+1}}\;\DTS[n^{a_{m+1}}]
\]
with input superscript $b_1=1$. Speedup Rules~0 and~1 are unchanged from \cite[Definition~4.1]{W13} (Rule~2 is two applications of Rule~1 with $x=0$ and $x$, giving the same line): for $0\le x\le a-1$,
\begin{itemize}
\item if $m=0$, $\DTS[n^a]$ goes to $(Q\,n^x)^{\max(x,1)}(Q'\,n^0)^1\DTS[n^{a-x}]$;
\item if $m>0$, the innermost block $(Q_m\,n^{a_m})^{b_{m+1}}$ becomes $(Q_m\,n^{\max(a_m,x)})^{\max(x,b_{m+1})}(Q'\,n^0)^{b_{m+1}}$ and the kernel becomes $a-x$.
\end{itemize}
The $\QBF_k$ slowdown removes $j$ innermost blocks ($1\le j\le\min(k,m)$) and replaces the kernel by
\[
c\cdot\max\{a_{m+1},\; b_{m-j+1},\; a_i, b_{i+1} \text{ for } m-j+1\le i\le m\},
\]
where $b_{m-j+1}$ is the superscript in front of the outermost removed block (the input length of the removed $\Sigma_j$/$\Pi_j$ subcomputation). This is the literal checker \lean{stepL} and \lean{checkL} in \lean{CertDefs.lean}, which uses exact rational arithmetic and requires normal form (no intermediate quantifier-free lines).

\subsection{The relaxed calculus \lean{QStep}}

For the upper limit, we drop guess exponents (which never exceed their after-superscripts in legal literal lines), replace the input floor $1$ by $0$, allow $0\le x\le a$, and allow every output exponent to be raised arbitrarily (weakening). A line in \lean{Calculus.lean} is a pair \lean{QLine} $= \mathrm{List}\;\mathbb{R} \times \mathbb{R}$ storing the stack of after-superscripts \emph{innermost first} together with the kernel exponent $a$. Write $\operatorname{front}([]):=0$, $\operatorname{front}(b::l):=b$, and $\operatorname{lmax}(l):=\max(\{0\}\cup l)$.

\begin{definition}[Relaxed step \lean{QStep}]\label{def:qstep}
Fix $c\in\mathbb{R}$ and $k\in\mathbb{N}$. A step \lean{QStep c k} is either:
\begin{enumerate}
\item \textbf{Speedup} (\lean{QStep.speed}): from $(\mathrm{st},a)$ with $0\le x\le a$, move to $(b_2::b_1::\mathrm{st.tail},\; a_2)$ whenever $\max(x,\operatorname{front}(\mathrm{st}))\le b_1$, $\operatorname{front}(\mathrm{st})\le b_2$, and $a-x\le a_2$.
\item \textbf{Slowdown} (\lean{QStep.pop}): from $(\mathrm{rem}\mathbin{++}\mathrm{rest},\; a)$ with $1\le |\mathrm{rem}|\le k$, move to $(\mathrm{rest},\; e)$ whenever $c\cdot\max(a,\;\operatorname{lmax}(\mathrm{rem}),\;\operatorname{front}(\mathrm{rest}))\le e$.
\end{enumerate}
\end{definition}

Note that when $\mathrm{st}=[]$, $\mathrm{st.tail}=[]$ and $\operatorname{front}(\mathrm{st})=0$, so \lean{QStep.speed} creates two blocks $(b_2, b_1)$ with $x\le b_1$ and $0\le b_2$ in a single step, matching Speedup Rule~0.

\begin{proposition}[Soundness of the literal checker; \lean{checkL\_sound}]\label{prop:sound}
If $c\in\mathbb{Q}_{\ge0}$ and \lean{checkL c k a0 ops = true}, then there exists $a'\in\mathbb{R}$ with $0\le a_0$, $a'<a_0$, and a \lean{QStep c k} path from $([],a_0)$ to $([],a')$.
\end{proposition}

\section{Threshold parameterisation and grouping potential}\label{sec:potential}

\subsection{The threshold parameters \lean{Par}}

Fix $k\ge1$. Set $q:=k/c$. The equation $p_k(c)=0$ is equivalent to the pair of identities
\[
q\cdot c = k, \qquad c - k = q(2-q), \qquad 0 < q < 1
\]
stored in the Lean structure \lean{Par}. Indeed, substituting $c=k+q(2-q)$ into $q c = k$ yields $k(1-q)=q^2(2-q)$, whose right-hand side minus left-hand side is $-k<0$ at $q=0$ and $1>0$ at $q=1$.

\begin{lemma}[Existence and monotonicity; \lean{exists\_par}, \lean{ge\_of\_poly}, \lean{par\_one\_eq\_r7}]
For every integer $k\ge1$ there exists $P:\lean{Par}$ with $P.k=k$, satisfying $p_k(P.c)=0$ and $k<P.c\le k+1$. Every $c>k$ with $p_k(c)\ge0$ satisfies $P.c\le c$. When $k=1$, $P.c=2\cos(\pi/7)$.
\end{lemma}

\subsection{The grouping potential}

Fix $P:\lean{Par}$ with $k=P.k$, $c=P.c$, $q=P.q$. Define the two auxiliary functions (\lean{Lam}, \lean{Lf}):
\[
\Lambda(s,e) := \min\!\bigl(e,\;\max(q(s+e),\; cs)\bigr), \qquad
L_r(s,e) := \max\!\left(s,\;\frac{\Lambda(s,e)}{r+1}\right) \quad (r\in\mathbb{N}).
\]
A \emph{valid grouping} $g=[h_1,\dots,h_t]$ of a stack $\mathrm{st}$ (\lean{Valid k g st}) partitions $\mathrm{st}$ from innermost to outermost into consecutive blocks of lengths $1\le h_i\le k$. Its value \lean{val P st g a} is defined recursively by $\lean{val}(\mathrm{st}, [], a) = a$ and
\begin{align*}
&\;\lean{val}(\mathrm{st}, h::g, a) := \\
&\quad \lean{val}\!\Bigl(\mathrm{st.drop}\;h,\; g,\; c\cdot\max\bigl(\operatorname{lmax}((\mathrm{st.take}\;h).\mathrm{tail}),\; \operatorname{front}(\mathrm{st.drop}\;h),\; L_{k-h}(\operatorname{front}(\mathrm{st}), a)\bigr)\Bigr).
\end{align*}
Intuitively, a valid grouping anticipates how the remaining blocks on the stack will eventually be popped in groups of size at most $k$; a group of size $h<k$ still has $r=k-h$ spare slots that can absorb future speedups before the group is popped, reducing the effective kernel by the factor $1/(r+1)$ inside $L_r$.

\begin{lemma}[Core inequalities; \lean{key}, \lean{lf\_S}, \lean{lf\_E}, \lean{top\_eq}]\label{lem:core}
For all $s,x,D,K\ge0$:
\begin{enumerate}
\item (\lean{key}) If $K\ge c\cdot\max(\max(x,s),\;\Lambda(s,D)/k)$, then $\Lambda(s, x+D)\le \Lambda(\max(x,s), K)$.
\item (\lean{lf\_S}) For every $r\in\mathbb{N}$, if $K\ge c\cdot\max(\max(x,s),\;L_{k-1}(s,D))$, then $L_r(s,x+D)\le L_r(\max(x,s),K)$.
\item (\lean{lf\_E}) For every $r\ge1$, $L_r(s,x+D)\le \max(\max(x,s),\; L_{r-1}(s,D))$.
\item (\lean{top\_eq}) For every $a\ge0$, $c\cdot L_{k-1}(0, a) = a$.
\end{enumerate}
\end{lemma}

\begin{theorem}[Step monotonicity; \lean{step\_pot}]\label{thm:step}
Let $L,L':\lean{QLine}$ with $L$ nonnegative (\lean{NN L}) and \lean{QStep P.c P.k L L'}. For every valid grouping $g'$ of $L'.1$, there exists a valid grouping $g$ of $L.1$ such that
\[
\lean{val}(P, L.1, g, L.2) \le \lean{val}(P, L'.1, g', L'.2).
\]
\end{theorem}

\begin{proof}[Proof idea]
For a slowdown popping $j\in\{1,\dots,k\}$ blocks from $L.1=\mathrm{rem}\mathbin{++}\mathrm{rest}$, prepend $j$ to $g'$: the outermost values match because $L_{k-j}(\operatorname{front}(\mathrm{rem}),a)\le\max(\operatorname{front}(\mathrm{rem}),a)\le\max(a,\operatorname{lmax}(\mathrm{rem}),\operatorname{front}(\mathrm{rest}))$. For a speedup pushing a new block onto a nonempty stack, inspect the innermost group size $h$ of $g'$: if $h=1$, the new block forms its own group and Lemma~\ref{lem:core}(\lean{lf\_S}) applies; if $h\ge2$, removing the new block shrinks the group to size $h-1$ and Lemma~\ref{lem:core}(\lean{lf\_E}) applies. When $L.1=[]$, the new stack has length $2$, and \lean{top\_eq} collapses the two possible groupings ($[1,1]$ or $[2]$ when $k\ge2$) directly to $a$.
\end{proof}

\section{Main theorems}\label{sec:main}

By induction along a path using Theorem~\ref{thm:step} and monotonicity of \lean{QStep} in $c$ (\lean{qstep\_restrict}), starting from $L_0=([],a_0)$ and ending at $L_n=([],a')$ (where the unique valid grouping is $[]$, with value $a'$), we obtain $a_0\le a'$.

\begin{theorem}[Upper limit for all $k\ge1$; \lean{qbf\_no\_refutation}, \lean{qbf\_no\_refutation\_strict}]\label{thm:limit}
Let $k\ge1$ be an integer and let $c\in\mathbb{R}$ with $c>k$.
\begin{enumerate}
\item If $c^3/k - c^2 - 2c + k \ge 0$, then every \lean{QStep c k} path from $([],a_0)$ ($a_0\ge0$) to $([],a')$ satisfies $a_0\le a'$.
\item If $c^3/k - c^2 - 2c + k > 0$, then every nonempty \lean{QStep c k} path ($n\ge1$) from $([],a_0)$ ($a_0>0$) to $([],a')$ satisfies $a_0 < a'$.
\end{enumerate}
\end{theorem}

For $k=1,2,3,4$, \lean{Cert.lean} encodes four literal refutations from Williams's recursive word family, verified in the Lean kernel with \lean{decide +kernel} (\lean{qcert1\_ok}--\lean{qcert4\_ok}). Combining Proposition~\ref{prop:sound}, downward monotonicity in $c$ (\lean{refutes\_mono}), and Theorem~\ref{thm:limit} gives the bracketed intervals in \lean{Summary.lean}:

\begin{theorem}[Certified threshold intervals; \lean{interval\_k1}--\lean{interval\_k4}, \lean{cert\_below}]
Define \lean{Refutes c k} to hold when there exist $0\le a_0$ and $a'<a_0$ with a \lean{QStep c k} path from $([],a_0)$ to $([],a')$. Then:
\begin{itemize}
\item $k=1$: \lean{Refutes c 1} holds for $0\le c\le 7/4$ and fails for $c\ge 2\cos(\pi/7)$;
\item $k=2$: \lean{Refutes c 2} holds for $0\le c\le 289/100$ and fails for $c>2$ with $c^3/2-c^2-2c+2\ge0$;
\item $k=3$: \lean{Refutes c 3} holds for $0\le c\le 393/100$ and fails for $c>3$ with $c^3/3-c^2-2c+3\ge0$;
\item $k=4$: \lean{Refutes c 4} holds for $0\le c\le 989/200$ and fails for $c>4$ with $c^3/4-c^2-2c+4\ge0$.
\end{itemize}
\end{theorem}

\section{Scope and limitations}\label{sec:scope}

\begin{itemize}
\item \textbf{Modelling claim.} Williams \cite[Section~5]{W13} states the $\QBF_k$ slowdown only briefly in prose. Our literal checker \lean{checkL} generalizes \cite[Definition~4.1]{W13} directly, retaining the in-front superscript $b_{m-j+1}$ (as in Williams's Table~I LP for $k=1$). Dropping guess exponents and input floors in \lean{QStep} only enlarges the set of proofs, so the impossibility theorem covers any tighter literal reading.
\item \textbf{Lower side for general $k$.} Williams's Theorem~5.1 gives a paper construction achieving every $c<c_k$ for all $k\ge1$; we kernel-check four explicit rational instances ($k=1,2,3,4$) rather than formalising the general parametric family.
\item \textbf{What is outside the formalisation.} Soundness of the complexity-theoretic speedup and slowdown lemmas, rules beyond Speedup Rules~0--2 and $k$-block Slowdown, and lower bounds obtained outside this calculus are not formalised.
\end{itemize}

\begin{thebibliography}{9}
\bibitem[BW15]{BW15}
S.~Buss and R.~Williams, \emph{Limits on alternation-trading proofs for time-space lower bounds}, Computational Complexity \textbf{24} (2015), 533--600.

\bibitem[W13]{W13}
R.~Williams, \emph{Alternation-trading proofs, linear programming, and lower bounds}, ACM Transactions on Computation Theory \textbf{5} (2013), no.~2, Article~6 (and CMU full version).
\end{thebibliography}

\end{document}
